% Generated draft: requires mathematical and source review.
\documentclass[11pt]{article}
\usepackage[margin=25mm]{geometry}
\usepackage{fontspec}
\setmainfont{Latin Modern Roman}
\usepackage{amsmath,amssymb,mathtools,unicode-math}
\setmathfont{Latin Modern Math}
\usepackage{enumitem,needspace,xurl,hyperref}
\setlength{\parindent}{0pt}
\setlength{\parskip}{5pt}
\setlength{\emergencystretch}{3em}
\newcommand{\N}{\mathbb{N}}
\newcommand{\Z}{\mathbb{Z}}
\newcommand{\Q}{\mathbb{Q}}
\newcommand{\R}{\mathbb{R}}
\newcommand{\C}{\mathbb{C}}
\newcommand{\F}{\mathbb{F}}
\newcommand{\A}{\mathbb{A}}
\newcommand{\PP}{\mathbb{P}}
\newcommand{\E}{\mathbb{E}}
\newcommand{\eps}{\varepsilon}
\newcommand{\dd}{\,\mathrm{d}}
\begin{document}
\section*{UnsolvedMath ID: 30000760}
\textbf{Source problem code:} catalogue entry
\subsection{Definitions and mathematical statement}
Let $G$ be reductive, let $G_{\mathrm{aff}}$ denote the corresponding affine Kac--Moody group intended by the source, and let $G_{\mathrm{aff}}^\vee$ denote its affine Langlands dual. For a specified tensor category $\operatorname{Rep}(G_{\mathrm{aff}}^\vee)$, its representation Grothendieck ring $K_0(\operatorname{Rep}(G_{\mathrm{aff}}^\vee))$ has generators $[V]$, relations $[V]=[V']+[V'']$ for short exact sequences $0\to V'\to V\to V''\to0$, and multiplication induced by tensor product. The source writes this ring as $K(\operatorname{Rep}(G_{\mathrm{aff}}^\vee))$. Its proposed affine Satake setting is schematically \[\mathcal S_{\mathrm{aff}}:\mathcal H_{\mathrm{aff}}\xrightarrow{\ \sim\ }\widehat{K_0(\operatorname{Rep}(G_{\mathrm{aff}}^\vee))\otimes_{\Z}\C},\] where the hat denotes the source's unspecified completion and $\mathcal H_{\mathrm{aff}}$ denotes what it calls the affine Hecke algebra. The question is whether this Satake isomorphism admits an explicit formula analogous to Lusztig's explicit $q$-analogue formula. The source supplies neither the analogue itself nor its coefficient/basis normalization, so no particular coefficient identity is asserted here. The algebra on the left must be clarified: the ordinary full Iwahori affine Hecke algebra is generally noncommutative, whereas a representation ring of this kind is commutative; a spherical algebra or another specified Satake algebra is therefore essential to the intended reading, rather than an automatic identification with every affine Hecke algebra. The construction of $G_{\mathrm{aff}}$, the dual root datum, the representation category, level conditions, completion, base field, and role of $q$ are not specified in the extract. The supplied 2026-08-21 triage attributes the proposed analogue to Braverman, describes its affine or double-affine Grassmannian program as conjectural in subsequent papers, and labels the general formula open.
\subsection{Short English statement}
Ask for an explicit Lusztig-type $q$-analogue formula describing the intended affine Satake map. The extract does not fix the relevant Hecke algebra, representation category, completion, or coefficient normalization, so these must be recovered before a particular formula can be stated.
\subsection{Sources}
\begin{itemize}\item Local UnsolvedMath registry, permanent ID 30000760, source problem catalogue entry, \emph{Explicit Formulas for the Affine Satake Isomorphism}, proposed year 2007. No individual proposer is supplied; the literature assessment calls the analogue Braverman's conjecture.\item Supplied report citation: \emph{Algebraic Groups} (2008), Oberwolfach Reports, DOI 10.4171/owr/2007/22, pp.\ 1191--1242: \url{https://doi.org/10.4171/owr/2007/22}.\item Supplied literature label: Braverman--Finkelberg, double affine Grassmannian I, \url{https://arxiv.org/abs/0711.2083}.\item Supplied literature label: Braverman--Finkelberg, double affine Grassmannian II, \url{https://arxiv.org/abs/0908.3390}.\end{itemize}
\end{document}
